% ==============================================
% About file: Format and style
% Description: list of packages, colors \& Style
% ==============================================

\usepackage{amsmath,amssymb}             % AMS Math
\usepackage[T1]{fontenc}
\usepackage{rotating,booktabs}
\usepackage{xcolor}
\usepackage{caption}
\usepackage{subcaption}
\usepackage{paralist}
\usepackage{ulem}
\usepackage{etoolbox}
\usepackage{graphicx}
\usepackage{float}
\usepackage{multicol}
\usepackage{microtype}
\usepackage{url}
\usepackage{hyperref}

\usepackage[left=1.0in,right=1.0in,top=1.0in,bottom=1.0in,includefoot,includehead,headheight=13.6pt]{geometry}
\renewcommand{\baselinestretch}{1.05}

\usepackage{titlesec}
% \titlespacing*{\chapter}{0pt}{-80pt}{10pt}

\renewcommand{\appendixname}{Apéndice}

\usepackage{colortbl}
\usepackage{algorithm}
\usepackage[noend]{algorithmic}
\usepackage{multirow}
%\usepackage{slashbox}

\usepackage{aecompl}
\usepackage{adjustbox}

% % Gantt Diagram
\usepackage{pgfgantt}

% % Nomenclatura
\usepackage{nomencl}
\newcommand{\nomunit}[1]{\hfill[#1]}
% \newcommand{\nomeqref}[1]{\hfill( #1 )}
% \newcommand{\nompageref}[1]{\hfill p.~#1}

% Links in pdf
\usepackage{color}
\definecolor{linkcol}{rgb}{0,0,0.4}
\definecolor{citecol}{rgb}{0.5,0,0}

% My pdf code
\usepackage{ifpdf}


% \let\showhyphens\relax

% %%%%%%%%%%%%%%%%%%%%%%%%%%


\definecolor{darkgreen}{RGB}{0,176,80}
\definecolor{lightgreen}{RGB}{169,208,142}
\definecolor{beige}{RGB}{242,242,242}
\definecolor{blue2}{RGB}{189,215,238}
\definecolor{barblue}{RGB}{153,220,254}
\definecolor{groupblue}{RGB}{51,102,254}
\definecolor{linkred}{RGB}{165,0,33}
\definecolor{groupgreen}{RGB}{0,240,0}
\definecolor{bargreen}{RGB}{0,220,0}
\definecolor{tabla1}{RGB}{255,217,102}
\definecolor{tabla2}{RGB}{189,215,238}

\renewcommand\sfdefault{phv}
% % \renewcommand\mddefault{mc}
% % %\renewcommand\bfdefault{bc}
% % \setganttlinklabel{s-s}{S -> S}
% % \setganttlinklabel{f-s}{F -> S}
% % \setganttlinklabel{f-f}{F -> F}
% % \sffamily

% \yyyymmdddate
% %%\today

%  %%%%%%%%%%%%%%%%%%%%%%%%%%%%%%%%%%%%%%%%%%%%%%%%%%%%%%%%%%%


% % Change this to change the some Meta Data in the *.pdf file
% % See hyperref documentation for information on those parameters

\hypersetup{
    colorlinks=true,
    citecolor=citecol,
    linkcolor=linkcol,
    urlcolor=linkcol,
    bookmarksopen=true,
    pdfauthor={Luis Felipe Guevara Gómez},
    pdftitle={Msc Thesis FG},
    pdfsubject={Trabajo de Grado | Ingeniería Electrónica | Javeriana Cali}
}

% definitions.
% -------------------

\setcounter{secnumdepth}{3}
\setcounter{tocdepth}{2}


\usepackage{rotating}  % Sideways of figures & tables
% \usepackage{natbib}  % Put References at the end of each chapter
% % Do -NOT- put 'sectionbib' option here.
% % Sectionbib option in 'natbib' will do.

% % \usepackage{txfonts}    % Public Times New Roman text & math font

% =========================
% FANCY HEADER CLEAN
% =========================
% % Fancy Header Style Options
\usepackage{fancyhdr}       % Fancy Header and Footer
\setlength{\headheight}{30pt}
\pagestyle{fancy}           % Sets fancy header and footer
\fancyhf{}

% =========================
% PÁGINAS PARES (even)
% =========================
\fancyhead[LE]{\bfseries\thepage} % número izquierda
\fancyhead[RE]{\bfseries\nouppercase{\leftmark}} % capítulo derecha

% =========================
% PÁGINAS IMPARES (odd)
% =========================
\fancyhead[LO]{\bfseries\nouppercase{\rightmark}} % sección izquierda
\fancyhead[RO]{\bfseries\thepage} % número derecha

\renewcommand{\chaptermark}[1]{%
  \markboth{Capítulo \thechapter.\ #1}{}}

\renewcommand{\sectionmark}[1]{%
  \markright{\thesection.\ #1}}

% \let\headruleORIG\headrule
% \renewcommand{\headrule}{\color{black} \headruleORIG}

\fancypagestyle{plain}{
  \fancyhf{}
  \renewcommand{\headrulewidth}{0pt}
}


\renewcommand{\headrulewidth}{1.0pt}

\arrayrulecolor{black}
\fancypagestyle{plain}{
    \fancyhf{}
    \renewcommand{\headrulewidth}{0pt}
}
% %%% Fancy Header ==========================


\makeatletter
\def\cleardoublepage{%
  \clearpage
  \if@twoside
    \ifodd\c@page\else
      \hbox{}
      \thispagestyle{empty} % ← CLAVE
      \newpage
    \fi
  \fi
}
\makeatother


\newenvironment{maxime}[1]
{
    \vspace*{0cm}
    \hfill
    \begin{minipage}{0.5\textwidth}%
    %\rule[0.5ex]{\textwidth}{0.1mm}\\%
    \hrulefill $\:$ {\bf #1}\\
    %\vspace*{-0.25cm}
    \it
}%
{%
    \hrulefill
    \vspace*{0.5cm}%
    \end{minipage}
}

\newenvironment{bulletList}%
{ \begin{list}%
	{$\bullet$}%
	{\setlength{\labelwidth}{25pt}%
	 \setlength{\leftmargin}{30pt}%
	 \setlength{\itemsep}{\parsep}}}%
{ \end{list} }

% Some useful commands and shortcut for maths:  partial derivative and stuff
\newcommand{\pd}[2]{\frac{\partial #1}{\partial #2}}
\def\abs{\operatorname{abs}}
\def\argmax{\operatornamewithlimits{arg\,max}}
\def\argmin{\operatornamewithlimits{arg\,min}}
\def\diag{\operatorname{Diag}}
\newcommand{\eqRef}[1]{(\ref{#1})}
\renewcommand{\epsilon}{\varepsilon}

% centered page environment
\newenvironment{vcenterpage}
{
    \newpage\vspace*{\fill} \fancyhf{} \renewcommand{\headrulewidth}{0pt}
}
{\vspace*{\fill}\par\pagebreak}

% ==============================================
% About: Contacts details University
% Description: Datos de los autores, directores,
% decano, jurados, etc.
% ==============================================


% ==========================================
% Formato extra PUJ
% ==========================================

%Fecha
\newcommand{\fecha}[1]{\def \fecha{#1}}

%Nombre del Trabajo
\newcommand{\titulo}[1]{\def \titulovar{#1}}

%Nombre de los directores del proyecto
\newcommand{\directorproyecto}[2]{\def \director{#1 #2}}
%\newcommand{\codirectorproyecto}[2]{\def \codirector{#1 #2}}

%Nombre del director de carrera
\newcommand{\directorcarrera}[2]{\def \directorIE{#1 #2}}
%Primer autor
\newcommand{\autorA}[2]{\def \authorA{#2}\def \authorACOD{#1}}
\newcommand{\autorB}[2]{\def \authorB{#2}\def \authorBCOD{#1}}

% TeK
%Nombre del decano de la FIC
\newcommand{\decanoFIC}[2]{\def \deanFIC{#1 #2}}
%Jurados
\newcommand{\juradoA}[2]{\def \juryA{#1 #2}}
\newcommand{\juradoB}[2]{\def \juryB{#1 #2}}


\usepackage{xcolor}
% \definecolor{linkred}{RGB}{120,0,40}



% % =========================
% % SPACING + FLOAT CONTROL
% % =========================
% \setlength{\lineskip}{1pt}
% \setlength{\normallineskip}{1pt}
% \setlength{\parskip}{0pt plus 1pt}

% \setcounter{topnumber}{2}
% \renewcommand{\topfraction}{0.7}
% \setcounter{bottomnumber}{1}
% \renewcommand{\bottomfraction}{0.3}
% \setcounter{totalnumber}{3}
% \renewcommand{\textfraction}{0.2}
% \renewcommand{\floatpagefraction}{0.5}

% \setcounter{dbltopnumber}{2}
% \renewcommand{\dbltopfraction}{0.7}
% \renewcommand{\dblfloatpagefraction}{0.5}
